Calculation of the angular distance on the orbit which the planet travels
during the transit. \hfill(6 points)\\
Since the orbit is circular,
\[ \sin\frac{\alpha}{2} = \frac{\pi t}{T} \]
where $t$ and $T$ are, respectively, the duration of the transit and the
orbital period of the planet. We can leave the angle as $\sin\frac{\alpha}{2}$
as this will be useful later; the numerical value of this
expressiob % sic: source misprint for "expression"
is about $0.112$.

Calculation of the velocity of the planet. \hfill(10 points)\\
The velocity $v$ at the beginning and end of the transit and the measured
difference in velocities $\Delta v$ form an isosceles triangle, with the angle
between the equal sides $= 2\alpha$. Thus
\[ 2v\cdot\sin\frac{\alpha}{2} = \Delta v \]
from which we get a velocity equal to about $134$~km/h.
% sic: source says "km/h"; the value is in km/s
\hfill(2 points)

Calculation of the radius of the orbit. \hfill(2 points)\\
Since we know the velocity on a circular orbit, this is trivial:
\[ 2\pi r = vT \Rightarrow r = \frac{vT}{2\pi} \]
Substituting we get about $6.4\cdot10^{6}$~km. \hfill(2 points)

Calculation of the radius of the star. \hfill(2 points)\\
With the velocity on the orbit and the transit duration, we get the radius of
the star:
\[ 2R = tv \Rightarrow R = \frac{tv}{2} = 7.5\cdot10^{5}\ \mathrm{km}. \]
The star is thus similar to the Sun in size. \hfill(2 points)

Calculation of the mass of the % sic: sentence truncated in source
star:

Given the orbital speed and radius, do this in the usual way:
\[ \frac{v^2}{r} = \frac{GM}{r^2} \]
\hfill(2 points)

thus
\[ M = \frac{v^2 r}{G} \]
The mass is therefore about $1{,}8\cdot10^{30}$~kg, so again similar to the
mass of the Sun. \hfill(2 points)
